\rowcolors{2}{white}{udcgray!25}
\small
\setlength{\tabcolsep}{3pt}
\begin{longtable}{ p{2cm}| p{3cm}| p{2cm}| p{5.5cm} }
  \caption{Casos de uso de gestión de avisos}
  \label{tab:cu-gestion-avisos} \\

  \rowcolor{udcpink!25}
  \textbf{Identificador} & \textbf{Caso de uso} & \textbf{Estado} & \textbf{Descripción} \\\hline
  \endfirsthead
  \multicolumn{4}{c}{\tablename\ \thetable{} -- {\small \textit{(vén da páxina anterior)}}} \\
  \rowcolor{udcpink!25}
  \textbf{Identificador} & \textbf{Caso de uso} & \textbf{Estado} & \textbf{Descripción} \\\hline
  \endhead
  \multicolumn{4}{c}{\dotfill{\small \textit{(continúa na páxina seguinte)}}\dotfill} \\
  \endfoot
  \endlastfoot

  CU-29 & Crear un aviso de emergencia & Actualizado & Cualquier usuario, autenticado o no, podrá crear un aviso sobre una emergencia desde la localización en la que se encuentra y, opcionalmente, añadir imágenes. \\
  CU-30 & Listar avisos & Actualizado & Un usuario registrado podrá visualizar los avisos creados. \\
  CU-31 & Validar o denegar aviso & Actualizado & Un coordinador podrá aceptar o rechazar los avisos recibidos antes de convertirlos en emergencias. \\
  CU-32 & Eliminar aviso & Actualizado & Un coordinador podrá eliminar avisos que no hayan sido validados. \\
\end{longtable}
